% =====================================================================
% คำนำ — บอกว่าคู่มือนี้เป็นของระบบอะไร ใครควรอ่าน และแบ่งเนื้อหาอย่างไร
% ไม่ควรเกิน 1 หน้า (ชื่อผู้จัดทำจะพิมพ์ท้ายหน้าให้อัตโนมัติ)
% =====================================================================
\begin{guide}
คำนำบอกว่าคู่มือเป็นของระบบอะไร รุ่นใด ผู้อ่านแต่ละกลุ่มควรอ่านบทใด และช่องทางแจ้งข้อผิดพลาด
ไม่เกิน 1 หน้า ไม่ต้องเล่าความเป็นมาของโครงงาน
\end{guide}
เอกสารคู่มือการใช้งานฉบับนี้จัดทำขึ้นเพื่ออธิบายวิธีใช้งานระบบจัดการเอกสารโครงงานคอมพิวเตอร์แบบออนไลน์
เวอร์ชัน \systemversion{} ซึ่งช่วยให้นักศึกษาและอาจารย์ที่ปรึกษาจัดการข้อเสนอโครงงาน
เขียนเอกสารด้วย \LaTeX{} และติดตามการส่งงานได้ในที่เดียว

คู่มือแบ่งเนื้อหาตามผู้อ่าน 3 กลุ่ม ผู้ใช้งานทั่วไปอ่านบทที่~\ref{ch:start} ถึงบทที่~\ref{ch:advisor}
ซึ่งอธิบายการใช้งานตามบทบาทของนักศึกษา อาจารย์ที่ปรึกษา และผู้ดูแลรายวิชา
ผู้ดูแลระบบอ่านบทที่~\ref{ch:install} ความต้องการของระบบและการติดตั้ง และบทที่~\ref{ch:admin} การดูแลระบบ
ส่วนนักพัฒนาที่รับช่วงต่ออ่านบทที่~\ref{ch:dev} คู่มือสำหรับนักพัฒนา ผู้อ่านทุกกลุ่มเริ่มจากบทที่~\ref{ch:overview}
ภาพรวมของระบบ และใช้บทที่~\ref{ch:trouble} เมื่อพบปัญหา

ผู้จัดทำหวังเป็นอย่างยิ่งว่าคู่มือฉบับนี้จะช่วยให้ผู้ใช้งานเริ่มใช้ระบบได้อย่างถูกต้องและรวดเร็ว
หากพบข้อผิดพลาดหรือมีข้อเสนอแนะ สามารถแจ้งผู้จัดทำได้ตามช่องทางในหัวข้อ~\ref{sec:contact}
